\section{Citometría de flujo}
\label{section:flow_cytometry}

La citometría de flujo es una técnica que permite analizar de manera rápida las propiedades físicas de partículas individuales ---generalmente células--- a medida que fluyen frente a uno o varios láseres, mientras se encuentran en suspensión en una solución buffer. Durante el paso de cada célula por el láser, se registra la luz esparcida. Debido a su alta velocidad de adquisición y capacidad para analizar miles de células por segundo, la citometría de flujo constituye una herramienta ampliamente utilizada en hematología, inmunología, biología molecular y diagnóstico clínico.

Cuando una célula atraviesa el haz láser, parte de la radiación incidente es esparcida debido al contraste entre el índice de refracción de la célula y el medio circundante. La distribución angular de la luz dispersada depende de diversos factores, entre ellos el tamaño, la forma, la orientación y la distribución espacial del índice de refracción de la célula. En la mayoría de los citómetros convencionales se registran dos señales de esparcimiento. El esparcimiento frontal (Forward Scatter, FSC), medida en un pequeño intervalo angular alrededor de la dirección de propagación del haz incidente. Tradicionalmente, esta señal se relaciona con el tamaño efectivo de la célula, aunque también depende de su forma y de su índice de refracción y el esparcimiento lateral (Side Scatter, SSC), medida aproximadamente a $90^{\circ}$ especto al haz incidente. Esta señal se asocia comúnmente con la complejidad interna o granularidad celular, ya que es sensible a las variaciones del índice de refracción presentes dentro de la célula. Es importante señalar que los detectores no registran la intensidad dispersada en un único ángulo, sino la luz colectada dentro de un intervalo angular determinado por la geometría de las lentes y las aperturas del sistema óptico.

Un citómetro de flujo está compuesto por tres sistemas principales: el sistema de fluidos, el sistema óptico y el sistema electrónico.
\begin{itemize}
	\item Sistema fluídico: El sistema de fluidos tiene como objetivo transportar las partículas hasta el punto de intersección con el haz láser. Para ello, la muestra celular se inyecta en una corriente de fluido de envoltura (sheath fluid), la cual la confina en una corriente estrecha mediante el principio de enfoque hidrodinámico. Este proceso mantiene el flujo de la muestra centrado dentro del fluido de envoltura y permite que las partículas atraviesen la región de iluminación una a una, evitando que dos células sean analizadas simultáneamente.
	
	En años recientes se han desarrollado citómetros con enfoque acústico, en los cuales ondas ultrasónicas complementan el enfoque hidrodinámico para mejorar la alineación de las células, incrementar la velocidad de adquisición y reducir el riesgo de obstrucción del sistema.
	
	\item Sistema óptico: El sistema óptico está formado por uno o varios láseres, lentes de colección  y detectores. Los láseres iluminan las partículas que atraviesan la región de análisis, mientras que las lentes colectan la luz esparcida. La señal de esparcimiento frontal suele detectarse mediante un fotodiodo, ya que la intensidad de la luz dispersada en esta dirección es considerablemente mayor.
	
	\item Sistema electrónico: El sistema electrónico convierte las señales ópticas registradas por los detectores en señales digitales que posteriormente son procesadas por el software del instrumento. 
\end{itemize}

Las mediciones conjuntas de FSC y SSC permiten representar cada célula mediante un punto en un diagrama bidimensional, facilitando la identificación y separación de diferentes poblaciones celulares. En el caso de los leucocitos, por ejemplo, estas dos señales son suficientes para distinguir las principales subpoblaciones presentes en sangre periférica. En la instrumentación convencional, las señales de dispersión se obtienen mediante detectores colocados en regiones angulares específicas, principalmente en la dirección frontal y alrededor de $90^{\circ}$. Sin embargo, la distribución angular de la luz dispersada contiene información adicional sobre la geometría y las propiedades ópticas de la célula. En consecuencia, la exploración de otros intervalos angulares de detección podría mejorar la capacidad para diferenciar células con distintas morfologías, motivando el análisis desarrollado en este trabajo.